\section{3  A Taxonomy of Agent Failure Modes}

A critical prerequisite to structured remediation is structured cause-finding. We propose a four-category taxonomy of agent failure modes relevant to customer-facing agents. These categories are not mutually exclusive, but they suggest distinct remediation pathways, which is the operationally important property for APIP design. Following Kirkpatrick's [7] insight that evaluation should be hierarchical, we note that each failure mode category corresponds to a different level of the agent's behavioral stack, and consequently to a different intervention tier.

Table 1: Agent failure mode taxonomy with Kirkpatrick level mapping and remediation tier.

Behavioral Drift is perhaps the most common failure mode in production deployments. A customer service agent trained or prompted against a product catalog from six months ago will increasingly produce incorrect citations as the catalog evolves. This is not a model capability failure---it is a retrieval or context freshness failure, and its remediation pathway is correspondingly different from fine-tuning. In Kirkpatrick's terms, this is a Level 4 failure: the agent's learned behavior has not changed, but the organizational results are degrading because the environment has shifted.

Input Brittleness is particularly pernicious because aggregate metrics mask it. An agent with 85\% task success overall may be failing catastrophically on the 8\% of requests that involve a particular billing scenario or a non-native-English speaker segment. Without disaggregated evaluation, the APIP trigger will not fire until overall performance has substantially degraded. This maps to a Kirkpatrick Level 2 failure: the agent lacks the capability (learning) to handle a specific input class, even though its general capability is adequate.

Alignment Creep covers the class of failures in which the agent's behavior drifts from the intended value envelope---not necessarily producing factually incorrect outputs, but producing outputs that are off-brand, subtly non-compliant, or misaligned with the organization's customer communication policies. Hendrycks et al. [6] characterize this class of failure as value misalignment, distinct from capability failure, and argue it requires distinct detection and remediation mechanisms. In Kirkpatrick's terms, this is a Level 3 failure: the agent's behavior has changed in ways that do not serve the organizational goals, even though its underlying capability has not degraded.

Tool Misuse is increasingly relevant as agentic systems are given access to action surfaces: ticket creation, refund processing, escalation routing, and external API calls. An agent that correctly identifies user intent but invokes the wrong tool, invokes a tool with incorrect parameters, or fails to invoke a necessary tool produces a qualitatively different class of failure than a purely conversational agent---one with potentially significant downstream consequences.

